


\section{Zadanie 8 (Oskar Prusik)}
Do N ponumerowanych listów wkładamy losowo N liczb. Jakie jest prawdopodobieństwo, że żadna liczba nie trafi do odpowiedniego listu?

\subsection*{Rozwiązanie}
Całe zadanie to jest dosłownie zasada włączeń i wyłączeń.

Niech $A_i$ - zdarzenie, że $i$-ta liczba trafiła do $i$-tego listu (dla $i = 1, 2, \ldots, N$).

Z zasady włączeń i wyłączeń:
\[
\left| \bigcup_{i=1}^N A_i \right| = \sum_{\substack{I \subseteq [N] \\ I \neq \emptyset}} (-1)^{|I| + 1} \left| \bigcap_{i \in I} A_i \right|
\]

\subsection*{Obliczenie mocy przekrojów}

Dla ustalonego $I \subseteq [N]$ z $|I| = k$:
\[
\left| \bigcap_{i \in I} A_i \right| = (N-k)!
\]
ponieważ jeśli $k$ liczb jest na swoich miejscach, to pozostałe $N-k$ liczb można rozmieścić dowolnie.

Liczba takich podzbiorów $I$ to $\binom{N}{k}$, więc:
\[
\sum_{\substack{I \subseteq [N] \\ |I| = k}} \left| \bigcap_{i \in I} A_i \right| = \binom{N}{k} (N-k)!
\]

\subsection*{Moc zdarzenia że przynajmniej jedna liczba trafiła dobrze}

\[
\left| \bigcup_{i=1}^N A_i \right| = \sum_{k=1}^N (-1)^{k+1} \binom{N}{k} (N-k)!
\]

\subsection*{Prawdopodobieństwo że przynajmniej jedna liczba trafiła dobrze}

\[
P\left( \bigcup_{i=1}^N A_i \right) = \frac{\left| \bigcup_{i=1}^N A_i \right|}{N!} 
= \sum_{k=1}^N (-1)^{k+1} \frac{ \binom{N}{k} (N-k)!}{N!}
 = \sum_{k=1}^N (-1)^{k+1} \frac{1}{k!}
\]

\subsection*{Prawdopodobieństwo że żadna liczba nie trafiła dobrze}

\[
P\left( \text{żadna liczba nie trafiła do odpowiedniego listu} \right) = 1 - P\left( \bigcup_{i=1}^N A_i \right)
\]

\[
P\left( \text{żadna liczba nie trafiła do odpowiedniego listu} \right) = 1 - \sum_{k=1}^N (-1)^{k+1} \frac{1}{k!} = \sum_{k=0}^N \frac{(-1)^k}{k!}
\]

\subsection*{Odpowiedź}
\[
\boxed{ \sum_{k=0}^N \frac{(-1)^k}{k!} }
\]
